% ==============================================================================
% CAPÍTULO 06 - RESULTADOS OBTENIDOS
% ==============================================================================

\chapter{Resultados Obtenidos}
\label{cap:resultados_obtenidos}
\justifying

\section{Implementación del Sistema}

Como parte del desarrollo del proyecto, se implementó una plataforma web funcional orientada a la gestión digital de quejas y denuncias para la DDP. La solución desarrollada integra componentes frontend, backend, bases de datos y servicios especializados para la administración centralizada de expedientes digitales.

Asimismo, se desarrollaron mecanismos de autenticación, control de acceso basado en roles, generación de folios, gestión documental y seguimiento de expedientes, permitiendo validar el flujo operativo principal del sistema.

La implementación permitió establecer una infraestructura modular basada en servicios independientes, facilitando la integración entre los diferentes componentes de la plataforma y mejorando la mantenibilidad del sistema.

Adicionalmente, se integró un componente inteligente de clasificación automática orientado al análisis preliminar de las narrativas textuales proporcionadas por los usuarios durante el registro de quejas.

\section{Módulos Desarrollados}

Durante el desarrollo del proyecto se implementaron distintos módulos funcionales enfocados en cubrir las necesidades operativas y administrativas identificadas durante el análisis del problema.

\begin{table}[H]
\centering
\caption{Módulos implementados dentro de la plataforma}
\label{tab:modulos_implementados}
\renewcommand{\arraystretch}{1.4}
\begin{tabularx}{\textwidth}{|p{5cm}|X|c|}
\hline
\textbf{Módulo} & \textbf{Descripción} & \textbf{Estado} \\
\hline

Portal del Quejoso &
Permite registrar quejas, consultar expedientes y cargar evidencias digitales. &
Implementado \\
\hline

Sistema de Autenticación &
Gestiona el inicio de sesión, recuperación de acceso y control de usuarios. &
Implementado \\
\hline

Gestión de Quejas &
Administra el registro, modificación y seguimiento de expedientes. &
Implementado \\
\hline

Dashboard Administrativo &
Proporciona herramientas de monitoreo y gestión institucional. &
Implementado \\
\hline

Clasificador Inteligente &
Realiza clasificación preliminar automática mediante técnicas de PLN. &
Implementado \\
\hline

Expediente Digital &
Centraliza información documental y evidencias asociadas al caso. &
Parcial \\
\hline

Sistema de Reportes &
Genera estadísticas e indicadores institucionales. &
Parcial \\
\hline

\end{tabularx}
\end{table}

\section{Resultados del Clasificador Inteligente}

Uno de los principales resultados obtenidos durante el desarrollo del proyecto corresponde a la implementación y evaluación de un clasificador inteligente orientado al análisis preliminar de quejas y denuncias institucionales.

Para el entrenamiento y evaluación de los modelos se utilizó un corpus conformado por casos reales anonimizados, permitiendo validar el comportamiento del sistema en escenarios cercanos al contexto operativo institucional.

El proceso de preparación de datos incluyó etapas de limpieza textual, normalización, eliminación de palabras vacías, stemming y vectorización mediante técnicas de representación textual orientadas al procesamiento de lenguaje natural.

Posteriormente, se evaluaron distintos algoritmos de aprendizaje automático con la finalidad de identificar el modelo con mejor desempeño para la clasificación preliminar de competencia institucional.

Los resultados obtenidos permitieron validar la viabilidad del uso de técnicas de Procesamiento de Lenguaje Natural y aprendizaje automático como apoyo al análisis inicial de expedientes.

\begin{table}[H]
\centering
\caption{Resultados generales del clasificador inteligente}
\label{tab:resultados_clasificador}
\renewcommand{\arraystretch}{1.4}
\begin{tabularx}{\textwidth}{|p{4cm}|c|c|c|}
\hline
\textbf{Modelo} & \textbf{Precisión} & \textbf{Recall} & \textbf{F1-Score} \\
\hline

Naive Bayes & 0.84 & 0.82 & 0.83 \\
\hline

Support Vector Machine & 0.89 & 0.88 & 0.88 \\
\hline

Random Forest & 0.85 & 0.84 & 0.84 \\
\hline

\end{tabularx}
\end{table}

Los resultados obtenidos muestran que el modelo basado en \textit{Support Vector Machine} presentó el mejor desempeño general, alcanzando valores superiores en precisión y capacidad de clasificación respecto a otros modelos evaluados.

\section{Pruebas Funcionales}

Con la finalidad de validar el correcto funcionamiento del sistema, se realizaron distintas pruebas funcionales orientadas a verificar la operación de los módulos principales de la plataforma.

Las pruebas contemplaron escenarios relacionados con autenticación, registro de quejas, generación de folios, carga de evidencias, seguimiento de expedientes y gestión administrativa.

\begin{table}[H]
\centering
\caption{Resultados de pruebas funcionales}
\label{tab:pruebas_funcionales}
\renewcommand{\arraystretch}{1.4}
\begin{tabularx}{\textwidth}{|p{6cm}|X|c|}
\hline
\textbf{Funcionalidad} & \textbf{Resultado Esperado} & \textbf{Estado} \\
\hline

Registro de quejas &
El sistema registra correctamente la información capturada por el usuario. &
Correcto \\
\hline

Generación de folios &
Se genera automáticamente un identificador único para cada expediente. &
Correcto \\
\hline

Inicio de sesión &
Validación correcta de credenciales y acceso basado en roles. &
Correcto \\
\hline

Carga de evidencias &
El sistema permite adjuntar archivos y documentos digitales. &
Correcto \\
\hline

Seguimiento de expedientes &
Consulta adecuada del estado actual de cada expediente. &
Correcto \\
\hline

Notificaciones &
El sistema muestra alertas y actualizaciones relevantes al usuario. &
Correcto \\
\hline

\end{tabularx}
\end{table}

Los resultados obtenidos durante las pruebas funcionales permitieron validar el comportamiento esperado de los principales componentes del sistema y verificar la integración correcta entre frontend, backend y base de datos.

\section{Resultados Operativos}

La implementación de la plataforma permitió validar mejoras importantes relacionadas con la organización, trazabilidad y automatización de los procesos institucionales asociados a la gestión de expedientes digitales.

Asimismo, la centralización de la información facilitó el acceso estructurado a documentos, evidencias y actuaciones relacionadas con cada caso, reduciendo inconsistencias operativas y mejorando la visibilidad del estado de los expedientes.

De igual manera, el uso de herramientas automatizadas permitió disminuir actividades repetitivas asociadas a validaciones manuales, clasificación preliminar y seguimiento administrativo.

Los resultados obtenidos demuestran el potencial de la solución para contribuir a la modernización tecnológica de los procesos internos de la DDP.

\section{Evidencias del Sistema}

